\documentstyle[%


righttag%


,11pt%


,amssymb%


,russian%               remove in Eng.


,izvru%                 Set izven% in Eng.


% ,izvru-ks             for brief communications


]{amsart}





% 





\newcommand{\rang}{\operatorname{rang}}


\newcommand{\tts}[1]{}





\theoremstyle{plain}


\theorembodyfont{\it}


\newtheorem{thmnonum}{\hskip\parindent Теорема}


\renewcommand{\thethmnonum}{}





\theoremstyle{definition}


\theorembodyfont{\rm}


\newtheorem{notenonum}{\hskip\parindent Замечание}


\renewcommand{\thenotenonum}{}





%\hoffset=-15mm%                  For Eng.


 


\setcounter{page}{57}


\god{1999}


\nomer{2~(441)} %                        Remove in Eng.


%\nomer{Vol.\,43, No.\,2}%               Set in Eng.


%\str{\pageref{begin}--\pageref{end}}%  Set in Eng.


\udk{517.93}





\author{\it Г.Г.\,Исламов}


% NB:   ^^ remove in Eng.


\title{Об одном свойстве мультипликаторов \\


линейных периодических систем}





\date{}


%\pagestyle{myheadings}%         Set in Eng.


\pagestyle{plain} %              Remove in Eng.


\begin{document}


%\thispagestyle{plain}%          Set in Eng.


\maketitle


\label{begin}








Пусть в линейной системе


\begin{equation}


\Dot x(t) = A(t)x(t) + u(t),


\end{equation}


где $A(t)$ есть $\omega$-периодическая $n\times n$-матрица с


комплекснозначными локально суммируемыми элементами, управление


$u(t)$ формируется по методу обратной связи


\begin{equation}


u(t) = - B(t)F(t)^{-1}x(t),


\end{equation}


где $F(t)=X(t)e^{-tK}$, $K=\frac{1}{\omega}\ln X(\omega)$


(здесь берется главное


значение логарифма), $X(t)$ --- матрицант невозмущенной системы


\begin{equation}


\Dot x(t) = A(t)x(t),


\end{equation}


а $B(t)$ есть $\omega$-периодическая $n\times n$-матрица с


комплекснозначными локально суммируемыми элементами.





При подстановке (2) в (1) получим однородную систему


\begin{equation}


\Dot x(t) = [A(t) - B(t)F(t)^{-1}]x(t)


\end{equation}


с $\omega$-периодической матрицей.


Свойства решений системы (4) будут определять свойства системы (1)


с обратной связью.





Пусть $E$ --- единичная матрица порядка $n$, $\Bbb C^n$ --- пространство


$n$-мерных векторов с комплексными компонентами.


Ранг функциональной матрицы $B(t)$\; $\rang B$ определим как


размерность образа отображения $(\Gamma z)(t) = B(t)z$, $z\in \Bbb C^n$,


$t\in [0,\omega]$.


Обозначим $\ker (X(\omega) - \rho E) = \{z\in \Bbb C^n :


X(\omega)z = \rho z\}$. По определению (напр., \cite{1}, c.\,28)


матрица $X(\omega)$ называется матрицей монодромии, а ее собственные


значения --- мультипликаторами системы (3). Совокупность


мультипликаторов называется спектром уравнения (3).





\begin{lem} Если возмущенная система $(4)$ не имеет


мультипликаторов в заданной области $\Omega$ комплексной плоскости


$\Bbb C$, то


\begin{equation}


\rang B \ge \max_{\rho\in \Omega}\dim\ker(X(\omega) - \rho E).


\end{equation}


\end{lem}





\begin{pf} Пусть $\rho\in\Omega$ и $z_1,\dotsc,z_m$


--- базис подпространства $\ker(X(\omega) - \rho E)$.


Допустим, что для некоторых скаляров $c_1,\dotsc,c_m$ при почти всех 


$t\in [0,\omega]$ линейная комбинация 


$\sum\limits_{k=1}^m c_kB(t)z_k = B(t)z =0$, где


$z=\sum\limits_{k=1}^m c_k z_k$. Отсюда для $x(t)=X(t)z=F(t)e^{tK}z=


\rho^{\frac{t}{\omega}}F(t)z$ имеем


$B(t)F(t)^{-1}x(t)=0$ и, значит, $\Dot x(t)=[A(t)-B(t)F(t)^{-1}]x(t)$.


Кроме того, $x(\omega)=X(\omega)z=\rho z=\rho x(0)$.


По условию леммы 1 в области $\Omega$ нет мультипликаторов системы (4).


Следовательно, $x(0)=z=0$ и в силу линейной независимости $z_1,\dotsc,z_m$


все $c_k=0, k=\overline{1,m}$. Это означает,


что матричный оператор $B(t)$ переводит линейно независимую систему векторов


$z_1,\dotsc,z_m$, $m=\dim\ker(X(\omega) - \rho E)$, в линейно независимую


на $[0,\omega]$ систему вектор-функций $B(t)z_1,\dotsc,B(t)z_m$. 


Отсюда следует (5).


\end{pf}





Следующая лемма аналогична лемме на с.\,302 монографии \cite{1}.





\begin{lem} Для любого комплексного $\lambda$


\begin{equation}


\ker(K-\lambda E) = \ker(X(\omega) - e^{\lambda\omega}E).


\end{equation}


\end{lem}





\begin{pf} Пусть $f(\lambda)=e^{\omega\lambda}$.


Так как производная $\Dot f(\lambda)$ не равна нулю на спектре матрицы $K$,


то в силу теоремы 9 (\cite{2}, с.\,158) при переходе от матрицы


$K$ к матрице $f(K)=e^{\omega K}=X(\omega)$ элементарные делители не


``расщепляются'',


т.\,е.~если матрица $K$ имеет элементарный делитель


$(\lambda-\lambda_j)^{m_j}$, то матрица монодромии $X(\omega)$ имеет


элементарный делитель $(\lambda-e^{\lambda_j\omega})^{m_j}$  и все


элементарные делители этой матрицы могут быть получены подобным образом.


Отсюда видно, что число $\alpha$ жордановых клеток в каноническом


представлении матрицы $K$, отвечающих $\lambda_j$, равно числу $\beta$


жордановых клеток в каноническом представлении матрицы $X(\omega)$,


отвечающих $e^{\lambda_j\omega}$ (напомним, что при определении матрицы


$K=\frac{1}{\omega}\ln X(\omega)$ мы берем главное значение логарифма и,


значит, $\lambda_k-\lambda_j\ne\frac{2\pi i}{\omega}l$ при


$\lambda_k\ne\lambda_j$, где $l$ целое). Так как $\alpha=


\dim\ker(K-\lambda_jE)$ и $\beta=\dim\ker(X(\omega)-e^{\lambda\omega}E)$,


то подпространства в (6) имеют одинаковую размерность.





Пусть $Kz=\lambda z$. Тогда $X(\omega)z=e^{\omega K}z=e^{\lambda\omega}z$.


Следовательно, $\ker(K-\lambda E)\subseteq \ker(X(\omega)-


e^{\lambda\omega}E)$. Так как размерности указанных подпространств


совпадают, то имеет место (6).


\end{pf}





\begin{lem} Пусть 


\begin{equation}


B(t)=F(t)S,


\end{equation}


где $S$ --- произвольная постоянная $n\times n$-матрица.


Тогда $Y(t)=F(t)e^{t(K-S)}$ будет матрицантом возмущенной


системы $(4)$ и $Y(\omega)=e^{\omega(K-S)}$ --- матрицей


монодромии этой системы.


\end{lem}





{\bf Доказательство.} Для $F(t)=X(t)e^{-tK}$ имеем


\begin{multline*}


\Dot F(t)=\Dot X(t)e^{-tK}-X(t)Ke^{-tK}=A(t)X(t)e^{-tK}-X(t)e^{-tK}K= \\


A(t)F(t)-F(t)K=[A(t)-B(t)F(t)^{-1}]F(t)-F(t)(K-S).


\end{multline*}


Так как $Y(t)=F(t)e^{t(K-S)}$, то с учетом найденного представления


для производной $\Dot F(t)$ получим


\begin{multline*}


\Dot Y(t)=\Dot F(t)e^{t(K-S)}+F(t)(K-S)e^{t(K-S)}= \\


=\{[A(t)-B(t)F(t)^{-1}]F(t)-F(t)(K-S)\}e^{t(K-S)}+F(t)(K-S)e^{t(K-S)}= \\


=[A(t)-B(t)F(t)^{-1}]F(t)e^{t(K-S)}=[A(t)-B(t)F(t)^{-1}]Y(t).


\quad\square


\end{multline*}





Следующее утверждение составляет содержание теоремы двойственности


заметки \cite{3} в комплексном случае.





\begin{lem} Пусть $\Lambda$ --- произвольное собственное


подмножество комплексной плоскости $\Bbb C$. Тогда


\begin{equation}


\max_{\lambda\in\Lambda}\dim\ker(K-\lambda E)=\min\rang S,


\end{equation}


где минимум берется по всем матрицам $S$ порядка $n$, для которых


в области $\Lambda$ нет собственных значений матрицы $K-S$.


\end{lem}





Основным результатом данной работы является





\begin{thmnonum} Пусть $\Omega$ --- произвольная область $\Bbb C$,


содержащая не все ненулевые комплексные числа. Тогда


$$


\max_{\rho\in\Omega}\dim\ker(X(\omega)-\rho E)=\min\rang B,


$$


где минимум берется по всем $\omega$-периодическим $n\times n$-матрицам


$B(t)$ с комплекснозначными локально суммируемыми компонентами, для которых


возмущенная система $(4)$ не имеет мультипликаторов в заданной области


$\Omega$.


\end{thmnonum}





\begin{pf} В силу леммы 1 достаточно показать, что в (5)


оценка снизу достигается. Выберем максимально широкое множество


$\Lambda$ так, что $\Omega=\{\rho: \rho=e^{\omega\lambda},


\lambda\in\Lambda\}$. В силу (6) и (8)


\begin{equation}


\max_{\rho\in\Omega}\dim\ker(X(\omega)-\rho E)=


\max_{\lambda\in\Lambda}\dim\ker(K-\lambda E)=\min\rang S,


\end{equation}


где минимум берется по всем матрицам $S$ порядка $n$, для которых в области


$\Lambda$ нет собственных значений матрицы $K-S$. Для матрицы минимального


ранга $S=S_0$ в (9) по формуле (7) построим


$\omega$-периодическую матрицу $B(t)$ такую, что для матрицы монодромии


возмущенной системы (4) будем иметь $Y(\omega)=e^{\omega(K-S_0)}$.


Так как в области $\Lambda$ нет собственных значений матрицы $K-S_0$,


то в области $\Omega$ не может быть собственных значений матрицы $Y(\omega)$.


В силу (7) $B(t)=F(t)S_0$. Отсюда видим, что


ранг функциональной матрицы $B(t)$ равен величине


$$


\rang S_0=\max_{\rho\in\Omega}\dim\ker(X(\omega)-\rho E).


$$


Это доказывает, что равенство в (5) достигается.


\end{pf}





\begin{notenonum} Пусть $Q(t)$ --- $\omega$-периодическая


$n\times n$-матрица с комплекснозначными локально суммируемыми


элементами. Отсутствие у возмущенной системы


$\Dot x(t) = [A(t) - Q(t)]x(t)$ мультипликаторов в заданной области


$\Omega$ эквивалентно отсутствию при любом $\rho\in\Omega$


ее нетривиальных решений, обладающих свойством


$x(t+\omega) = \rho x(t)$. При выполнении одного из этих эквивалентных


условий функциональная матрица $Q(t)F(t)$ в силу леммы~1 должна иметь


ранг, не меньший максимальной геометрической кратности мультипликаторов


из области $\Omega$ невозмущенной системы $\Dot x(t) = A(t)x(t)$.


Теорема показывает, что минимально возможное значение ранга совпадает с


максимальной геометрической кратностью.


\end{notenonum}





\begin{thebibliography}{9}





\bibitem{1}


Якубович В.А., Старжинский В.М. {\em Параметрический резонанс в линейных


системах}. -- М.: Наука, 1987. -- 328 с.


\bibitem{2}


Гантмахер Ф.Р. {\em Теория матриц}. -- 3-е изд. -- М.: Наука, 1967. -- 576 с.


\bibitem{3}


Исламов Г.Г. {\em Об управлении спектром динамической системы} // Дифференц.


уравнения. -- 1987. -- Т.\,23. -- \No\,8. -- C.\,1299--1302.





\end{thebibliography}





\vskip 2cm





\post{\it Удмуртский государственный}{\it Поступила}


\post{\it университет}{25.06.1997}





\label{end}


\end{document}


